\documentclass[12pt, a4paper]{exam}
\usepackage{fontspec}
\usepackage{xunicode}
\usepackage{xltxtra}
\usepackage{polyglossia}
\usepackage{graphicx}
\usepackage{tikz}
\usetikzlibrary{positioning, calc}
\usepackage{geometry}
\usepackage{enumitem}
\usepackage{multicol}
\usepackage{array}
\usepackage{lastpage}
\usepackage{amssymb}
\usepackage{multirow}
\usepackage{amsmath}

\setmainlanguage{thai}
\XeTeXlinebreaklocale "th"
\XeTeXlinebreakskip = 0pt plus 1pt
\linespread{1.25}

% Page Margins
\geometry{top=0.8in, bottom=1in, left=0.8in, right=0.8in}

% Fonts
\newfontfamily\thaifont[
  Scale=1.33,
  Script=Thai,
  Path = ../fonts/,
  Extension = .ttf,
  BoldFont = *-Bold,
  ItalicFont = *-Italic,
  BoldItalicFont = *-BoldItalic
]{THSarabunNew}

\setmainfont[
  Scale=1.33,
  Path = ../fonts/,
  Extension = .ttf,
  BoldFont = *-Bold,
  ItalicFont = *-Italic,
  BoldItalicFont = *-BoldItalic
]{THSarabunNew}

% Header and Footer
\pagestyle{headandfoot}
\extraheadheight{0.2in}

\newcommand{\examheader}{%
    ชื่อ-นามสกุล \makebox[3.0in]{\hrulefill} รหัสนักศึกษา \makebox[2.0in]{\hrulefill}
}

\firstpageheader{\examheader}{}{}
\runningheader{\examheader}{}{}

\firstpagefooter{\textit{ข้อสอบรายวิชา MTH234 Linear Algebra}}{}{\thepage}
\runningfooter{\textit{ข้อสอบรายวิชา MTH234 Linear Algebra}}{}{\thepage}

% Custom Thai Choice Labels
\makeatletter
\newcommand{\thaialph}[1]{%
  \ifcase\value{#1}\or ก\or ข\or ค\or ง\or จ\else\@ctrerr\fi
}
\makeatother

% Question Format
\renewcommand{\thequestion}{\arabic{question}}
\renewcommand{\questionlabel}{\thequestion.}

% Choice Format
\renewcommand{\thechoice}{\thaialph{choice}}
\renewcommand{\choicelabel}{\thechoice)}

\renewcommand{\choiceshook}{%
  \setlength{\leftmargin}{2.5em}%
  \setlength{\labelwidth}{20pt}%
  \setlength{\labelsep}{5pt}%
}

\begin{document}
\sloppy
\setcounter{page}{2}

\noindent \textbf{\underline{คำสั่ง}} : ให้นักศึกษาทำข้อสอบทุกข้อ โดยทำลงไปในข้อสอบ
\vspace{0.5cm}

\begin{questions}

% ข้อ 1 (Matrix Operations & Algebra)
\question \textbf{(5 คะแนน)} กำหนดให้
\[
A = \begin{bmatrix} -2 & 0 & 4 \\ 2 & 6 & 3 \\ 4 & -1 & 2 \end{bmatrix}, \quad
B = \begin{bmatrix} 1 & 0 & -3 \\ -2 & -5 & 0 \\ -2 & 1 & 2 \end{bmatrix}, \quad \text{และ} \quad
C = \begin{bmatrix} 0 & -2 & 0 \\ -1 & -3 & 4 \end{bmatrix}
\]
จงหา:
\begin{parts}
    \part \textbf{(1 คะแนน)} \mbox{เมทริกซ์} $3A - B^T$ (โดยที่ $T$ คือ ทรานสโพส (transpose))
    \vspace{4cm}
    \part \textbf{(2 คะแนน)} \mbox{เมทริกซ์} $BC$ และ $CB$ (ถ้าสามารถหาค่าได้)
    \vspace{6cm}
    \part \textbf{(2 คะแนน)} \mbox{เมทริกซ์} $X$ ที่ทำให้ $X(A+B) = C$
\end{parts}
\vspace{\stretch{1}}

\newpage

% ข้อ 2 (Row Reduction & RREF)
\question \textbf{(5 คะแนน)} กำหนดให้
\[
A = \begin{bmatrix} 
1 & 0 & -2 & 1 & 0 \\ 
0 & -1 & -3 & 1 & 3 \\ 
-2 & -1 & 1 & -1 & 3 \\ 
0 & 3 & 9 & 0 & -12 
\end{bmatrix}
\]
\begin{parts}
    \part \textbf{(4 คะแนน)} {\addfontfeature{LetterSpace=-1.5} จงใช้ลำดับของการดำเนินการตามแถวเบื้องต้น (elementary row operations) เพื่อหาเมทริกซ์} ที่เป็นรูปแบบขั้นบันไดตามแถวแบบลดรูป (reduced row echelon form) ซึ่งสมมูลกันทางแถว (row equivalent) กับเมทริกซ์ $A$
    \vspace{12cm}
    \part \textbf{(1 คะแนน)} จงระบุแรงค์ (rank) ของ\mbox{เมทริกซ์} $A$
\end{parts}
\vspace{\stretch{1}}

\newpage

% ข้อ 3 (Inverse & Determinant)
\question \textbf{(10 คะแนน)} กำหนดให้
\[
A = \begin{bmatrix} 
2 & 6 & 6 & 2 \\ 
2 & 7 & 3 & 6 \\ 
1 & 5 & 0 & 1 \\ 
3 & 7 & 0 & 7 
\end{bmatrix}
\]
จงใช้การดำเนินการตามแถวเบื้องต้น (elementary row operations) เพื่อหา:
\begin{parts}
    \part \textbf{(4 คะแนน)} ค่าดีเทอร์มิแนนต์ของ $A$ $(\operatorname{\text{\fontfamily{cmr}\selectfont det}}(A))$
    \vspace{\stretch{1}} 
    \newpage
    \part \textbf{(6 คะแนน)} \mbox{เมทริกซ์}ผกผันของ $A$ ($A^{-1}$) (ถ้ามี)
\end{parts}
\vspace{\stretch{1}}

\newpage

% ข้อ 4 (System of Linear Equations)
\question \textbf{(8 คะแนน)} พิจารณาระบบสมการเชิงเส้น:
\[
\begin{aligned}
2x_1 + 3x_2 + 5x_3 &= 4 \\
3x_1 + 5x_2 + 9x_3 &= 7 \\
5x_1 + 9x_2 + 13x_3 &= 17
\end{aligned}
\]
จงหาผลเฉลยของระบบสมการนี้โดย:
\begin{parts}
    \part \textbf{(4 คะแนน)} การใช้ลำดับของการดำเนินการตามแถวเบื้องต้น (using elementary row operations)
    \vspace{\stretch{1}}
    \newpage
    \part \textbf{(4 คะแนน)} การใช้กฎของคราเมอร์ (using Cramer's rule)
\end{parts}
\vspace{\stretch{1}}

\newpage

% ข้อ 5 (Homogeneous System & Basis)
\question \textbf{(6 คะแนน)} กำหนดให้
\[
A = \begin{bmatrix} 
1 & 2 & 3 \\ 
2 & -1 & 1 \\ 
3 & 1 & 4 \\ 
5 & 1 & 6 
\end{bmatrix}, \quad
\vec{b} = \begin{bmatrix} 
0 \\ 0 \\ 0 \\ 0 
\end{bmatrix}
\]
จงตรวจสอบว่าระบบสมการ $A\vec{x} = \vec{b}$ มีผลเฉลยที่ไม่ชัด (non-trivial solutions) \mbox{หรือไม่} พร้อมทั้งหา:
\begin{parts}
    \part \textbf{(3 คะแนน)} ผลเฉลยทั่วไป (general solution)
    \vspace{\stretch{1}}
    \newpage
    \part \textbf{(3 คะแนน)} ฐานหลัก (basis) สำหรับปริภูมิผลเฉลย (solution space) และระบุมิติ (dimension)
\end{parts}
\vspace{\stretch{1}}

\newpage

% ข้อ 6 (Parameters & Consistency)
\question \textbf{(6 คะแนน)} จงหาค่าของ $k$ ทั้งหมดที่ทำให้ระบบสมการเชิงเส้นต่อไปนี้:
\[
\begin{aligned}
x + y + kz &= 3 \\
x + ky + z &= 2 \\
kx + y + z &= 1
\end{aligned}
\]
เป็นไปตามกรณีต่อไปนี้:
\begin{parts}
    \part \textbf{(2 คะแนน)} ไม่มีผลเฉลย (no solution)
    \vspace{\stretch{1}}
    \part \textbf{(2 คะแนน)} มีผลเฉลยชุดเดียว (unique solution)
    \vspace{\stretch{1}}
    \part \textbf{(2 คะแนน)} มีผลเฉลยมากมายไม่จำกัด (infinitely many solutions)
\end{parts}
\vspace{\stretch{1}}

\end{questions}
\end{document}





